\uuid{fuJX}
\exo7id{1708}
\titre{exo7 1708}
\auteur{barraud}
\organisation{exo7}
\datecreate{2003-09-01}
\isIndication{false}
\isCorrection{false}
\chapitre{Réduction d'endomorphisme, polynôme annulateur}
\sousChapitre{Applications}
\cpge{mpsi-mp2i.nombres-reels-suites.i, pcsi.nombres-reels-suites.g, ptsi.nombres-reels-suites.g, mpsi-mp2i.nombres-complexes.e, pcsi.nombres-complexes.e, ptsi.nombres-complexes.e, mpsi-mp2i.polynomes-fractions-rationnelles.c, pcsi.polynomes.c, ptsi.polynomes.c, bcpst1.suites-reelles-usuelles, mp-mpi.reduction-endomorphismes.d, pc.reduction.c, psi.reduction-endomorphismes.c, pt.reduction-endomorphismes.c, mp-mpi.reduction-endomorphismes.c, pc.reduction.b, psi.reduction-endomorphismes.b, pt.reduction-endomorphismes.b, bcpst2.valeurs-propres-vecteurs-propres.b}
\module{Algèbre}
\niveau{L2}
\difficulte{}

\contenu{
\question{
Déterminer toutes les suites $(u_{n})_{n\in\N}$ à valeur complexes
  telles que~:
  $$
  \forall n\in\N, u_{n+3}+2 u_{n+2}+2u_{n+1}+u_{n}=0.
  $$
  Montrer que les suites réelles satisfaisant cette relation sont les
  suites de la forme~:
  $$
  u_{n}=A(-1)^{n}+B\cos(\frac{2n\pi}{3}+\phi)
  $$
  où $A,B$ et $\phi$ sont des réels.
}
}
